% =====================================================================
% Secao de Implementation (nova). Descreve o sistema concreto, fundamentada
% no codigo do repositorio. Idioma: ingles. Sugestao: inserir logo apos a
% Secao "DNAT Architecture" e antes da "Security Analysis".
% =====================================================================

\section{Implementation}
\label{sec:implementation}

This section describes how the three planes are realized, so that the security properties
evaluated later can be traced to concrete mechanisms rather than to design intent alone. The
system is deployed as three containers on a shared runtime---\texttt{dnat-client} (CVM1),
\texttt{dnat-builder} (CVM2), and \texttt{dnat-executor} (CVM3)---each standing in for the
hardware-backed confidential VM~\cite{amd2020sevsnp} of a production deployment, and is
contrasted with two baselines: a single-container deployment that colocates the same assets
(A) and a no-isolation variant that removes the microVM entirely (M0). The consequences of
that substrate substitution are examined in Section~\ref{sec:evaluation}.

\subsection{Control Plane (CVM1)}

The control plane hosts the marketplace API, a local IPFS node, and the blockchain client,
exposing registration, access checks, and execution orchestration through a small HTTP API.
On registration, CVM1 does not install dependencies itself: it assembles a build bundle (the
source plus a manifest of declared dependencies) and forwards it to the builder. When the
builder returns an \texttt{application.ext4} artifact, CVM1 encrypts it with AES-256-GCM
over a gzip-compressed image, in a versioned envelope carrying the IV, authentication tag,
and metadata. Only the encrypted blob and a manifest are published to IPFS; the plaintext
artifact never leaves the control plane. Access control is enforced on chain: an execution
request first calls \texttt{hasAccessByIds}, and the workflow aborts with \emph{access
denied} if the caller has not purchased access to both assets. A purchased right is recorded
under the key \texttt{keccak256(encryptedDatasetHash, encryptedApplicationHash,
buyer)}---bound to the encrypted IPFS CIDs of the two assets and the buyer's address rather
than to mutable identifiers, so the grant is stable under off-chain re-indexing---and is
verified before any bundle is assembled.

\subsection{Build Plane (CVM2)}

The builder exposes a single \texttt{/build} endpoint that accepts a bundle and launches an
ephemeral Firecracker microVM to resolve dependencies. Unlike the executor, the build
microVM is given a network interface, because resolving and compiling Python wheels may
require egress to package indexes. Egress is constrained on the host side: a TAP device is
created and iptables rules \texttt{REJECT} traffic from the guest toward private ranges
(\texttt{10.0.0.0/8}, \texttt{172.16.0.0/12}, \texttt{192.168.0.0/16},
\texttt{127.0.0.0/8}, \texttt{169.254.0.0/16}, and the host TAP address), while the
remaining outbound traffic is masqueraded. Inside the guest, an artifact builder produces a
read-only \texttt{ext4} image holding the application source, its \texttt{site-packages}, a
lock file of resolved dependencies, and a marker file used later for disk discovery.
Dependencies are first wheel-built and then installed with \texttt{--no-index} from the
local wheelhouse; reusable wheels are persisted in a size-pruned cache that is the builder's
only durable state---and, as Section~\ref{sec:evaluation} notes, its only cross-request
channel. After the build, the host returns only the artifact and its metadata; the
per-request build state is discarded.

\subsection{Execution Plane (CVM3)}

The executor exposes a single \texttt{/execute} endpoint. For each request it copies a fresh
root overlay, creates new input and output \texttt{ext4} disks, optionally attaches the
application artifact read-only, and launches a Firecracker microVM with one vCPU, 512~MiB of
memory and---crucially---no network interface. The guest does not receive disk roles out of
band; it discovers them by scanning candidate block devices for marker files, then extracts
the workspace and runs \texttt{workspace/run.sh}, whose output is persisted to the output
disk. Host and guest synchronize over the serial console: the guest writes an
\texttt{EXECUTION\_COMPLETE} marker, after which the host issues a graceful shutdown,
escalating to termination if needed, then mounts the output disk and derives the HTTP
response solely from the persisted result. A \texttt{cleanup} trap unmounts every loop device
and removes the per-run working directory on exit, which is what makes execution state
ephemeral.

\subsection{Application Contract, Asset Integrity, and Baselines}

An application is registered as unmodified source together with a manifest of its
third-party dependencies, which the builder resolves and bundles into the read-only
\texttt{application.ext4} artifact. Because dependencies are shipped as ordinary
\texttt{site-packages} inside that artifact---rather than linked into a trusted-execution
enclave at build time---an application that follows this contract runs with arbitrary
declared dependencies, unchanged across all three environments and without recompilation or
enclave-specific adaptation (R6). This is the reuse property that
Section~\ref{sec:evaluation} contrasts with the enclave-based original DNAT.

Assets are bound to the chain by content hash: registration records the SHA-256 of the
artifact, which lets a consumer detect tampering between IPFS and the contract. The first
baseline (environment~A) runs the identical control, build, and execution logic, including
Firecracker isolation for both microVMs, but inside a single container whose filesystem
holds the IPFS repository, the wheel cache, and the executions store, and whose environment
holds the asset-encryption key and the wallet private key. Its only structural difference
from the compartmentalized deployment is the colocation of these assets---exactly the
variable the security evaluation isolates.

The second baseline removes isolation altogether. With \texttt{DNAT\_NO\_MICROVM} set, the
executor runs the same \texttt{workspace/run.sh} contract directly in the container's
namespace---no Firecracker, no per-run disks, no network removal---while emitting the same
result, so the request path is otherwise unchanged. In this variant (M0) the application
inherits the container's network, environment (including the control-plane secrets), and
persistent filesystem, which is what lets the evaluation separate the microVM's operational
cost from the containment it provides.
